%% ----------------------------------------------------------------
%% Thesis.tex
%% ---------------------------------------------------------------- 
\documentclass[sotoncolour]{uosthesis}      % Use the Thesis Style with custom link colour
\usepackage{lmodern}
\usepackage{physics}
\usepackage{cleveref}
\usepackage{graphicx}
\usepackage{multirow}
\usepackage{svg}

\usepackage{nomencl,etoolbox,ragged2e}
\graphicspath{{resources/}}   % Location of your graphics files
\usepackage[round]{natbib}            % Use Natbib style for the refs.
\usepackage{bibentry}          % Use bibentry for prepublished works
\nobibliography*               % Use bibentry for prepublished works
\usepackage{attrib}            % Use the attrib package for quotations
\hypersetup{colorlinks=true}   % Set to false for black/white printing
\usepackage[version=4]{mhchem}

% -----------------------------------------------------------------
% Custom legend line markers for figure captions (examples, edit or remove)

\newcommand{\solidgreenline}{\raisebox{2pt}{\tikz{\draw[-,black!20!green,solid,line width = 2.0pt](0,0) -- (5mm,0);}}}

\newcommand{\solidblueline}{\raisebox{2pt}{\tikz{\draw[-,black!5!blue,solid,line width = 2.0pt](0,0) -- (5mm,0);}}}

\newcommand{\solidredline}{\raisebox{2pt}{\tikz{\draw[-,black!5!red,solid,line width = 2.0pt](0,0) -- (5mm,0);}}}

\newcommand{\dottedredline}{\raisebox{2pt}{\tikz{\draw[-,black!5!red,dotted,line width = 2.0pt](0,0) -- (5mm,0);}}}

\newcommand{\dashedredcircleline}{\raisebox{2pt}{\tikz{\draw[black!100!red, dashed, line width=0.5pt] (0,0) -- (5mm,0); \node[draw, circle, inner sep=1pt, fill=black!100!red] at (2.5mm, 0) {};}}}

\newcommand{\dashedredopencircleline}{\raisebox{2pt}{\tikz{\draw[black!100!red, dashed, line width=0.5pt] (0,0) -- (5mm,0); \node[draw, circle, inner sep=1pt, fill=white!100!red] at (2.5mm, 0) {};}}}

\newcommand{\dashedblueplusline}{\raisebox{-0.5ex}{\tikz{\draw[dashed, black!100!blue, line width=0.5pt] (0,0) -- (5mm,0) node[midway] {\texttt{+}};}}}

\newcommand{\dashedbluestarline}{\raisebox{-0.5ex}{\tikz{\draw[dashed, black!100!blue, line width=0.5pt] (0,0) -- (5mm,0) node[midway] {\textasteriskcentered};}}}


\input{definitions}            % Include your abbreviations
%% ----------------------------------------------------------------
%% --------------------THESIS/DOC INFORMATION ---------------------
\department  {[School Name]}
\DEPARTMENT  {\MakeUppercase{\deptname}}
\group       {[Research Group Name]}
\GROUP       {\MakeUppercase{\groupname}}
\faculty     {[Faculty Name]}
\FACULTY     {\MakeUppercase{\facname}}
\title      {[Thesis Title]}
%% TODO: Replace with your name removing []
\authors    {[Author Name]} % Use of Soton Email unadvised, use ORCiD instead.
\addresses  {\groupname\\\deptname\\\univname}
\date       {\today}
%% Optional Fields TODO: Replace these fields with your own data

% \qualifications{[Insert Previous Qualifications, for example MEng]}
\orcidid{0000-0000-0000-0000}
\doi{10.0000/00000000.xxx}
\volume{i}{i} %Optional Volume Numbering Volume n of m
\subject    {}
\keywords   {}

%% ------------------------------ glossary addition
\usepackage{longtable,floatrow,booktabs}
% \usepackage{hyperref}
\usepackage{siunitx}

\usepackage[nogroupskip,toc,indexonlyfirst]{glossaries-extra}
\usepackage{ifthen}

\newboolean{printglossaryflag}
\setboolean{printglossaryflag}{true} % Set the flag to true to print the glossary

%=============================================================================
%   G L O S S A R Y   S E T U P
%=============================================================================


\glsaddkey
{unit}          % new key
{\relax}        % default value if "unit" isn't used in \newglossaryentry
{\glsentryunit} % analogous to \glsentrytext
{\Glsentryunit} % analogous to \Glsentrytext
{\glsunit}      % analogous to \glstext
{\Glsunit}      % analogous to \Glstext
{\GLSunit}      % analogous to \GLStext

\glsaddkey
{def}       % new key
{\relax}            % default value if "def" isn't used in \newglossaryentry
{\glsentrydef}      % analogous to \glsentrytext
{\Glsentrydef}      % analogous to \Glsentrytext
{\glsdef}           % analogous to \glstext
{\Glsdef}           % analogous to \Glstext
{\GLSdef}           % analogous to \GLStext

\newglossary[alg]{acronym}{acr}{acn}{List of Acronyms}
\newglossary[nlgG]{notationGreek}{notG}{ntnG}{Greek Symbols}
\newglossary[nlgR]{notationRoman}{notR}{ntnR}{Roman Symbols}
\newglossary[dnlg]{dimensionlessNumber}{dnt}{dntn}{Dimensionless Numbers}
\newglossary[sslg]{superscript}{sst}{sstn}{Superscripts}
\newglossary[sblg]{subscript}{sbt}{sbtn}{Subscripts}

% If compiling manually (rather than via latexmk or Overleaf), each glossary
% needs its own makeindex call, e.g.:
% makeindex -s main.ist -t main.glg  -o main.gls  main.glo
% makeindex -s main.ist -t main.alg  -o main.acr  main.acn
% makeindex -s main.ist -t main.nlgG -o main.notG main.ntnG
% makeindex -s main.ist -t main.nlgR -o main.notR main.ntnR
% makeindex -s main.ist -t main.dnlg -o main.dnt  main.dntn
% makeindex -s main.ist -t main.sslg -o main.sst  main.sstn

% Useful references for the glossaries setup:
% http://tex.stackexchange.com/questions/78016
% http://tex.stackexchange.com/questions/71549
% http://tex.stackexchange.com/questions/8946
% http://tex.stackexchange.com/questions/26563
% http://tex.stackexchange.com/questions/126323

% Ensure table captions stay at the top
\floatsetup[table]{capposition=top, objectset=raggedright,
    margins=raggedright, midcode=captionskip, captionskip=10pt}

\newglossarystyle{customacronyms}
{
    \setglossarystyle{long4colheader}%
    \renewcommand*{\glossaryheader}{%
        \textbf{Symbol} & \textbf{Description}
        \tabularnewline
        \endhead
        \endfoot
    }%
    \setlength\glsdescwidth{.4\textwidth}%
    \setlength\glspagelistwidth{.95\textwidth}%
    \setlength{\extrarowheight}{3pt}%
    \renewcommand{\glossentry}[2]{\bfseries\glsentryitem{##1}\glstarget{##1}{\glossentryname{##1}} & \glossentrydesc{##1}
        \tabularnewline
    }%
}

\newglossarystyle{customsymbols}
{
    \setglossarystyle{long3colheader}%
    \renewcommand*{\glossaryheader}{%
        \textbf{Symbol} & \textbf{Description} & \textbf{Units}
        \tabularnewline
        \endhead
        \endfoot
    }%
    \setlength\glsdescwidth{.65\textwidth}%
    \setlength\glspagelistwidth{.95\textwidth}%
    \setlength{\extrarowheight}{3pt}%
    \renewcommand{\glossentry}[2]{\glsentryitem{##1}\glstarget{##1}{\glossentryname{##1}} & \glossentrydesc{##1} & \glsentryunit{##1}
        \tabularnewline
    }%
}

\newglossarystyle{customdimlessnum}
{
    \setglossarystyle{long3colheader}%
    \renewcommand*{\glossaryheader}{%
        \textbf{Symbol} & \textbf{Description} & \textbf{Definition}
        \tabularnewline
        \endhead
        \endfoot
    }%
    \setlength\glsdescwidth{.5\textwidth}%
    \setlength\glspagelistwidth{25\textwidth}%
    \setlength{\extrarowheight}{20pt}%
    \renewcommand{\glossentry}[2]{\glsentryitem{##1}\glstarget{##1}{\glossentryname{##1}} & \glossentrydesc{##1} & \glsentrydef{##1}
        \tabularnewline
    }%
}

\makeglossaries

\begin{document}
%% ------------------ FRONT MATTER ORGANISATION -------------------
\pagenumbering{gobble} % removes page number
\copyrightDeclaration{} % !!! Comment this line when printing the hardcopy !!!
% \raggedright                  %% Left justification; remove for full justification
                              %% Must be done after copyrightDeclaration
\frontmatter
\maketitle
\begin{abstract}

% Abstract text goes here.

\end{abstract}
\tableofcontents
\listoffigures
\listoftables

%% ---------- AUTHORSHIP DECLARATION / ACKNOW. / DEDICATORY ----------
%% Either include citations like below (as many as required, spaced with commas or 'and').
%% \bibentry command must be used here with prepublished papers
% \authorshipdeclaration{\bibentry{key1}\newline\bibentry{key2}}
%% Or state no citations:
% \authorshipdeclaration{}
%% -----------------------
\acknowledgements{
% Acknowledgements text goes here.
}
\dedicatory{[Dedication]}

%% ------------------ GLOSSARIES --------------------
%====================================================================
% Acronyms

\newglossaryentry{NASA}{
  type={acronym},
  sort={NASA},
  name={NASA},
  short={NASA},
  long={National Aeronautics and Space Administration},
  first={National Aeronautics and Space Administration (NASA)},
  description={United States civilian space agency}
}

\newglossaryentry{DNS}{
  type={acronym},
  sort={DNS},
  name={DNS},
  short={DNS},
  long={Direct Numerical Simulation},
  first={Direct Numerical Simulation (DNS)},
  description={Numerical solution resolving all relevant flow scales}
}

%====================================================================
% Greek symbols

\newglossaryentry{gamma}{
  type={notationGreek},
  name={$\gamma$},
  description={Ratio of specific heats},
  first={Ratio of specific heats ($\gamma$)},
  unit={1}
}

\newglossaryentry{density}{
  type={notationGreek},
  name={$\rho$},
  description={Fluid density},
  first={Fluid density ($\rho$)},
  unit={$kg\,m^{-3}$}
}

\newglossaryentry{viscosity}{
  type={notationGreek},
  name={$\mu$},
  description={Dynamic viscosity},
  first={Dynamic viscosity ($\mu$)},
  unit={$N\,s\,m^{-2}$}
}

%====================================================================
% Roman symbols

\newglossaryentry{velocity}{
  type={notationRoman},
  name={$u$},
  description={Velocity},
  first={Velocity ($u$)},
  unit={$m\,s^{-1}$}
}

\newglossaryentry{pressure}{
  type={notationRoman},
  name={$p$},
  description={Pressure},
  first={Pressure ($p$)},
  unit={$Pa$}
}

\newglossaryentry{temperature}{
  type={notationRoman},
  name={$T$},
  description={Temperature},
  first={Temperature ($T$)},
  unit={$K$}
}

%====================================================================
% Subscripts and superscripts

\newglossaryentry{freestream}{
  type={subscript},
  name={$\infty$},
  description={Free-stream condition}
}

\newglossaryentry{species}{
  type={subscript},
  name={$s$},
  description={Species index}
}

\newglossaryentry{dimensional}{
  type={superscript},
  name={$*$},
  description={Dimensional quantity}
}

%====================================================================
% Dimensionless numbers

\newglossaryentry{Re}{
  type={dimensionlessNumber},
  name={Re},
  description={Reynolds number},
  def={$\rho u L / \mu$}
}

\newglossaryentry{Pr}{
  type={dimensionlessNumber},
  name={Pr},
  description={Prandtl number},
  def={$\mu c_p / \kappa$}
}

%====================================================================


\chapter*{Definitions and abbreviations}
\addcontentsline{toc}{chapter}{Definitions and Abbreviations}

\glsaddall
\printglossary[type=acronym,style=customacronyms]
\printglossary[type=notationGreek,style=customsymbols]
\printglossary[type=notationRoman,style=customsymbols]
\printglossary[type=superscript,style=customacronyms]
\printglossary[type=subscript,style=customacronyms]
\printglossary[type=dimensionlessNumber,style=customdimlessnum]

\mainmatter
%% ------------------ MAIN MATTER (CONTENT) --------------------

\include{source/2.introduction/introduction}
% \include{source/3.methods/methods}
% \include{source/4.results/results}
%% ----------------------------------------------------------------
%% Conclusions.tex
%% ---------------------------------------------------------------- 
% \chapter{Conclusions and Future Work}\label{Chapter:Conclusions}

% \section{Conclusions}

\chapter{Conclusions and Future Work}\label{Chapter:Conclusions}

\section{Conclusions}

\section{Future Work}


% ----------------------------------------------------------------
% appendix
\appendix
%% ----------------------------------------------------------------
%% AppendixA.tex
%% ---------------------------------------------------------------- 
\chapter{Appendix A}
\label{Appendix:AppendixA}


\chapter{Appendix B}
\label{Appendix:AppendixB}


\backmatter

\bibliographystyle{apalike}
\bibliography{UOS}
\end{document}
%% ----------------------------------------------------------------